\section{The golden operator and its cubic shadows}
\label{sec:golden-operator}

The incidence cover supplies a golden orientation bit.  Relative to the
marked bridge datum, we attach that bit to the six-axis carrier and apply
three elementary operations to its
conference operator: middle exterior power, commutator with the diagonal
algebra, and centered squaring.  These operations produce the Joubert,
Segre, Igusa, and diagonal Clebsch models in one calculation.
The main route first fixes the outer family and identifies four exact
descriptions of its marked cubic.  A reader following the
source--shadow--return argument may then continue directly to
Section~\ref{sec:harmonic}.  The later exchange-spectrum and reconstruction
subsections prove independent structural consequences of the same conference
carrier.

Let \(X\) be the ordered six-axis set and
\(A_X=\Q^X/\Q\mathbf1\).  The outer six-set \(\mathcal T\) consists of the
six synthematic totals, that is, the six labelled one-factorizations of the
fixed labelled graph \(K_X\).
These are the classical mystic pentagons.  For each \(T\in\mathcal T\), its
complementary triangle colouring assigns signs
\(\epsilon_T(S)\in\{\pm1\}\) to the three-subsets \(S\subset X\), with the
global sign fixed by the marked outer orientation.  The signs obey the
four-point two-graph identity and transform by
\[
 \epsilon_{gT}(gS)=\operatorname{sgn}(g)\epsilon_T(S)
 \qquad(g\in S_X).
\]
Writing \(\epsilon_T(S)=(-1)^{\tau_T(S)}\), the four-point identity is
\begin{equation}
 \label{eq:two-graph-parity}
 \sum_{S\in\binom Q3}\tau_T(S)=0
 \qquad\left(Q\in\binom X4\right).
\end{equation}
This is the two-graph parity equation.  Switching a graph representative
leaves \(\tau_T\) unchanged; complementing \(\tau_T\) reverses every triangle
sign while preserving \eqref{eq:two-graph-parity}.
This is the signed coloured-triangle construction of the outer six-family
\cite[Sections~1.1--1.2 and~1.6]{HMSVOuter}.
The order-six conference two-graph is likewise classical: its switching
class has the pentagon-plus-isolated-vertex form and automorphism group
\(A_5\) \cite{BMS}.
The four-point identity reconstructs from each \(\epsilon_T\) a switching
class of symmetric zero-diagonal sign matrices; it is the unique order-six
conference class.  Choose representatives \(C_T\) whose triangle products
are \(\epsilon_T\), and transport the Hodge orientations with the ordered
axis basis.  Then \(C_T^2=5I\).  These choices constitute the coherent outer
marking.  The incidence sheet supplies only its orientation bit, not the
outer labels or representatives.

Put
\[
 J_T(x)=\sum_{|S|=3}\epsilon_T(S)x_S,
 \qquad x_S=\prod_{i\in S}x_i.
\]
The following table fixes every sign.  Take \(T_0\) to be the total attached
to the matrix displayed in Section~\ref{sec:orientation-source}, let
\(T_r=p_rT_0\), and write permutations in one-line notation.  In the column
headed ``coefficient word'', the signs occur in the order
\[
012,013,014,015,023,024,025,034,035,045,
123,124,125,134,135,145,234,235,245,345.
\]
\begin{equation}
\label{tab:outer-coefficient-words}
\begingroup
\newcommand{\sgminus}{\mathord{-}}
\newcommand{\sgplus}{\mathord{+}}
\begin{array}{c|c|c}
r&p_r&\text{coefficient word of }J_{T_r}\\ \hline
0&012345&\sgminus\sgminus\sgplus\sgplus\sgplus\sgminus\sgplus\sgplus\sgminus\sgminus\sgplus\sgplus\sgminus\sgminus\sgplus\sgminus\sgminus\sgminus\sgplus\sgplus\\
1&012354&\sgplus\sgplus\sgminus\sgminus\sgminus\sgminus\sgplus\sgplus\sgminus\sgplus\sgminus\sgplus\sgminus\sgminus\sgplus\sgplus\sgplus\sgplus\sgminus\sgminus\\
2&012435&\sgplus\sgminus\sgplus\sgminus\sgplus\sgminus\sgminus\sgminus\sgplus\sgplus\sgminus\sgminus\sgplus\sgplus\sgplus\sgminus\sgplus\sgminus\sgplus\sgminus\\
3&012453&\sgminus\sgplus\sgminus\sgplus\sgplus\sgplus\sgminus\sgminus\sgminus\sgplus\sgminus\sgplus\sgplus\sgplus\sgminus\sgminus\sgminus\sgplus\sgminus\sgplus\\
4&012534&\sgminus\sgplus\sgplus\sgminus\sgminus\sgplus\sgplus\sgminus\sgplus\sgminus\sgplus\sgminus\sgplus\sgminus\sgminus\sgplus\sgplus\sgminus\sgminus\sgplus\\
5&012543&\sgplus\sgminus\sgminus\sgplus\sgminus\sgplus\sgminus\sgplus\sgplus\sgminus\sgplus\sgminus\sgminus\sgplus\sgminus\sgplus\sgminus\sgplus\sgplus\sgminus
\end{array}
\endgroup
\end{equation}
Indeed,
\(\epsilon_{pT_0}(S)=\operatorname{sgn}(p)
\epsilon_{T_0}(p^{-1}S)\); applying this to the twenty triangle products of
the displayed matrix gives the six rows.  They are the signed
coloured-triangle Joubert frame of
\cite[Sections~1.1--1.2 and~1.6]{HMSVOuter}.  Thus this citation supplies no
unstated scalar or sign convention: its six coordinates are exactly the
forms \(J_T\) in \eqref{tab:outer-coefficient-words}.

\subsection{Four cubic shadows}
\label{sec:four-shadows}

For a three-subset \(S\subset X\), set
\[
 K_T=*\!\bigwedge\nolimits^3 C_T,
 \qquad
 Z_T(x)=\frac14\sum_{|S|=3}(K_T)_{SS}x_S,
 \qquad D_x=\operatorname{diag}(x).
\]
The Hodge star is taken in the ordered axis basis.  Writing \(X=\{0,\dots,5\}\)
in its order and \(S^{c}\) for the complementary triple in increasing order,
\[
 *e_S=\operatorname{sgn}(S,S^{c})\,e_{S^{c}},
 \qquad
 \operatorname{sgn}(S,S^{c})=(-1)^{\sigma(S)-3},
 \qquad
 \sigma(S)=\sum_{i\in S}i,
\]
where \(\operatorname{sgn}(S,S^{c})\) is the sign of the permutation listing
\(S\) and then \(S^{c}\).  The exponent counts that listing's inversions:
writing \(S=\{a_0<a_1<a_2\}\), the element \(a_k\) has \(k\) elements of \(S\)
below it, hence \(a_k-k\) of \(S^{c}\), and these sum to \(\sigma(S)-3\).
Since \(\sigma(S)+\sigma(S^{c})=15\), complementary signs multiply to
\((-1)^{9}=-1\),
which is \(*^2=-1\) in middle degree.  The orientation therefore travels with
the ordered basis: an odd relabelling or switching transports both \(K_T\) and
the sign of its diagonal, and it is not legitimate to switch the matrix while
holding the old Hodge convention fixed.  Figure~\ref{fig:four-cubic-shadows}
separates the four normalized formulas from their classical consequences.

\begin{figure}[t]
\centering
\begin{tikzpicture}[
  box/.style={draw,rounded corners=2pt,align=center,inner sep=4pt,
    font=\small,text width=39mm},
  center/.style={box,text width=43mm,thick},
  target/.style={box,text width=42mm},
  arr/.style={-{Latex[length=2mm]},semithick},
  eqline/.style={semithick},
  lab/.style={font=\scriptsize,fill=white,inner sep=1.2pt}
]
\node[box] (tri) at (-3.4,2.3) {triangle holonomy\\
  \(\sum_S\epsilon_T(S)x_S=Z_T\)};
\node[box] (wedge) at (3.4,2.3) {middle exterior power\\
  \(\frac14\sum_S(K_T)_{SS}x_S=Z_T\)};
\node[box] (pf) at (-3.4,0.5) {commutator Pfaffian\\
  \(\frac14\operatorname{Pf}[D_x,C_T]=Z_T\)};
\node[box] (det) at (3.4,0.5) {oriented cross-golden determinant\\
  \(10\sqrt5\det B_T(x)=Z_T\)};
\node[center] (z) at (0,-1.15) {four normalized formulas agree\\
  \(Z_T=J_T\)};
\draw[eqline,rounded corners=3pt] (tri.west) -- ++(-4mm,0) |- (z.north west);
\draw[eqline,rounded corners=3pt] (wedge.east) -- ++(4mm,0) |- (z.north east);
\draw[eqline] (pf.south east) -- (z.west);
\draw[eqline] (det.south west) -- (z.east);

\node[target] (outer) at (0,-2.75) {six outer coordinates\\
  \((Z_T)_{T\in\mathcal T}\)};
\node[target] (segre) at (-3.4,-4.45) {Segre cubic\\
  \(\sum_TZ_T=\sum_TZ_T^3=0\)};
\node[target] (igusa) at (3.4,-4.45) {centered squares \(W_T\)\\
  Segre--Igusa polar map};
\node[target] (clebsch) at (-3.4,-6.15) {hyperplane section \(z_T=0\)\\
  diagonal Clebsch cubic};
\draw[arr] (z) -- (outer);
\draw[arr] (outer.south west) -- node[lab,above left] {Segre relations}
  (segre.north);
\draw[arr] (outer.south east) -- node[lab,above right] {center}
  (igusa.north);
\draw[arr] (segre) -- node[lab,right] {section} (clebsch);
\end{tikzpicture}
\caption{The four constructions of the same marked cubic and their classical
targets.  Plain connectors denote the normalized identities displayed in the
upper boxes; arrows denote subsequent constructions.  The determinant-line
orientation is normalized by the cross-golden formula.}
\label{fig:four-cubic-shadows}
\end{figure}

\smallskip
\noindent\emph{What the four descriptions do and do not assert.}\par
For every symmetric matrix \(A\) of order six, the coefficient of \(x_S\) in
\(\operatorname{Pf}[D_x,A]\) is the signed complementary minor
\(\det A[S^c,S]\), hence the corresponding diagonal entry of
\(*\!\bigwedge^3A\).  The cross-golden determinant is the spectral-block form
of this identity once \(C^2=5I\).  Only the triangle comparison is
matrix-specific.
Theorem~\ref{thm:triangle-pfaffian-recognition} forces scalar-square rigidity
on the nondegenerate real locus and gives exact recognition on the sign locus.

For a six-tuple \((a_U)_{U\in\mathcal T}\), write
\[
 \operatorname{center}_T a_T
 =a_T-\frac16\sum_{U\in\mathcal T}a_U.
\]

\begin{theorem}[Golden operator and classical cubic shadows]
\label{thm:operator-shadows}
For a marked bridge datum and its coherent outer family:
\begin{enumerate}
\item The triangle products of \(C_T\) satisfy
\[
 (K_T)_{SS}=4(C_T)_{ij}(C_T)_{jk}(C_T)_{ki}=4\epsilon_T(S)
 \quad(S=\{i,j,k\}),
\]
so \(Z_T=J_T\) is the switching-invariant orientation cubic on \(A_X\).
The four-point two-graph identity and pair balance recover the conference
switching class; after fixing a vertex, the negative edges on the other five
vertices form a pentagon.
\item The same cubic is a commutator Pfaffian:
\[
 \operatorname{Pf}[D_x,C_T]=4Z_T(x),
 \qquad
 \det[D_x,C_T]=16Z_T(x)^2.
\]
Over \(E=\Q(\sqrt5)\), put
\(P_{T,\pm}=(I\pm C_T/\sqrt5)/2\) and
\(V_{T,\pm}=\operatorname{im}P_{T,\pm}\).  Define the restricted map
\[
 B_T(x)=P_{T,-}D_x\big|_{V_{T,+}}:V_{T,+}\longrightarrow V_{T,-}.
\]
Its basis-free determinant lies in
\(\det(V_{T,-})\otimes\det(V_{T,+})^{-1}\).  At \(T_0\), orient these
determinant lines so that the following equality has positive sign, and
transport the orientations through the coherent outer marking.  The inherited
symmetric forms and these orientations trivialize the determinant line, and
\[
 Z_T(x)=10\sqrt5\det B_T(x).
\]
Galois exchange swaps the two rank-three summands, so the determinant-line
contraction differs from the ordinary field norm by \((-1)^{3\cdot3}=-1\):
\[
 \det[D_x,C_T]
 =-8000\,N_{E/\Q}\!\left(\det B_T(x)\right).
\]
\item The six coordinates \((Z_T)_{T\in\mathcal T}\) satisfy
\[
 \sum_TZ_T=0,\qquad \sum_TZ_T^3=0.
\]
They are the signed Joubert coordinates and hence define the classical map
from six labelled points to the Segre cubic
\[
 \left\{[z]\in\mathbf P^5:\sum_Tz_T=\sum_Tz_T^3=0\right\}.
\]
Inside its hyperplane \(z_T=0\), the remaining linear relation cuts the
ambient space to \(\mathbf P^3\), and the remaining cubic equation is the
diagonal Clebsch cubic surface.
\item The centered determinants
\[
 W_T=Z_T^2-\frac16\sum_UZ_U^2
     =\frac1{16}\operatorname{center}_T\det[D_x,C_T]
\]
satisfy
\[
 \sum_TW_T=0,
 \qquad
 \left(\sum_TW_T^2\right)^2=4\sum_TW_T^4.
\]
Thus \(W\) is the Segre--Igusa polar map.
\end{enumerate}
\end{theorem}

\begin{proof}
The tight-frame argument in Section~\ref{sec:orientation-source} gives
\(C_T^2=5I\).  Its off-diagonal entries are the pair-balance identities
\[
 \sum_{k\ne i,j}(C_T)_{ij}(C_T)_{jk}(C_T)_{ki}=0.
\]
Triangle products are unchanged by diagonal switching and obey the
four-point identity because every edge of a four-set occurs twice.  In the
gauge \((C_T)_{0i}=1\), pair balance says that every vertex among the other
five has two negative incident edges.  The negative graph is therefore a
five-cycle.  This directly recovers the classical switching-class normal
form without an enumeration.

Expand the Pfaffian of \([D_x,C_T]\) over perfect matchings.  The coefficient
of \(x_S\) pairs \(S\) bijectively with \(S^c\), and its signed sum is the
complementary \(3\)-minor of \(C_T\).  The Hodge convention displayed above
makes that signed minor equal to \((K_T)_{SS}\).  There are two triangle
orbits, distinguished by \(\epsilon_T(S)=\pm1\); switching, transported
relabelling, and complementation reduce them to \(S=012\) and \(S=014\) in
the displayed matrix.  For these representatives the complementary blocks
\(C_T[S^c,S]\) are
\[
 \begin{pmatrix}1&-1&1\\-1&-1&1\\-1&-1&-1\end{pmatrix},
 \qquad
 \begin{pmatrix}1&-1&1\\1&-1&-1\\-1&-1&-1\end{pmatrix}.
\]
Their determinants are \(4\) and \(-4\); the Hodge signs
\(\operatorname{sgn}(S^c,S)\) are both \(-1\), while the internal triangle
products are \(-1\) and \(+1\).  Hence in both orbits
\((K_T)_{SS}=4(C_T)_{ij}(C_T)_{jk}(C_T)_{ki}\).  This exposes the
complementary-minor step rather than treating it as formal.  The Pfaffian is
therefore \(4Z_T\); squaring gives the determinant identity.

The projectors split the axis space into two three-spaces.  Relative to that
splitting,
\[
 [D_x,C_T]=
 \begin{pmatrix}
 0&-2\sqrt5\,B_T(x)^{\mathsf T}\\
 2\sqrt5\,B_T(x)&0
 \end{pmatrix}.
\]
Comparison with \(16Z_T^2\) gives
\(Z_T^2=500\det(B_T)^2\); the chosen determinant-line orientation fixes the
displayed sign.

The matching expansion gives \(Z_T=J_T\) coefficient by coefficient.
The table in \eqref{tab:outer-coefficient-words} therefore identifies the six operator cubics with the signed
Joubert coordinates, including their normalization and the rule
\(gZ_T=\operatorname{sgn}(g)Z_{gT}\).  Their Segre relations are classical
\cite[Sections~2.1--2.2]{HMSVOuter}, so
\(\sum Z_T=\sum Z_T^3=0\).  Setting one coordinate to zero gives
\(\sum_{U\ne T}Z_U=\sum_{U\ne T}Z_U^3=0\), the diagonal Clebsch cubic.
The centered-square formula and the symmetric Igusa equation are the
classical Segre--Igusa duality map \cite[Section~2.2]{HMSVOuter}, following
van der Geer.  Here the determinant identity in the theorem identifies that
polar vector with the centered commutator determinants.
\end{proof}

\begin{theorem}[Triangle--Pfaffian recognition]
\label{thm:triangle-pfaffian-recognition}
Let \(A=(a_{ij})\) be real symmetric of even order \(n\geq4\), with zero
diagonal and every off-diagonal entry nonzero.  Put
\[
 \mathcal T_A(x)=\sum_{|S|=3}
 \left(\prod_{\{i,j\}\subset S}a_{ij}\right)x_S,
 \qquad
 \Phi_A(x)=\operatorname{Pf}[D_x,A],
 \qquad x_S=\prod_{i\in S}x_i.
\]
If \(\Phi_A=\mu\mathcal T_A\) for some \(\mu\ne0\), then \(n=6\) and
\(A^2=\lambda I\) for some \(\lambda>0\).

If the off-diagonal entries also have a common absolute value, then, up to scale, switching,
and relabelling, \(A\) is the pentagon conference matrix; conversely
\(\Phi_A=\pm4\mathcal T_A\); relative to an orientation of the signed label space,
the sign is that space's two-valued orientation character.
\end{theorem}

\begin{proof}
Equality of degrees forces \(n/2=3\).  Since
\([D_{x+t\mathbf1},A]=[D_x,A]\), proportionality makes \(\mathcal T_A\)
translation-invariant.  Its coefficient of \(tx_ix_j\), for \(i\ne j\), is
\[
 a_{ij}\sum_{r\ne i,j}a_{ir}a_{rj}=a_{ij}(A^2)_{ij}.
\]
Thus \(A^2\) is diagonal.  Now \([A,A^2]=0\) and \(a_{ij}\ne0\) make its
diagonal constant.  Hence \(A^2=\lambda I\), where
\(\lambda=(A^2)_{ii}=\sum_{r\ne i}a_{ir}^2>0\).

In the equal-absolute-value case, rescale to a sign matrix \(B\), so
\(B^2=5I\).  The pentagon gauge argument at the start of the preceding proof
gives the class.  Its complementary-minor calculation gives the factor \(4\);
a signed relabelling multiplies it by that relabelling's determinant.
\end{proof}

The hypotheses exclude zero proportionality and vanishing off-diagonal entries.
What the theorem leaves open is not the conclusion \(A^2=\lambda I\), whose
solutions are described below, but which weighted matrices among them satisfy
the proportionality; those are not classified.

Both halves of the theorem have classical ingredients.  Which forms are
Pfaffians of skew matrices of linear forms is settled in general by Beauville
\cite{BeauvilleDeterminantal}, with the cubic-surface case constructive in
Tanturri \cite{TanturriPfaffian}, so a \(6\times6\) linear Pfaffian
representation of a cubic surface is a known object.  The triangle cubic's
coefficients are the order-three cycle products, half the principal
three-by-three minors of a hollow symmetric matrix
(\(\det A[S,S]=2a_{ij}a_{jk}a_{ki}\) for \(S=\{i,j,k\}\)), and cycle-sum
coordinates on principal minors are Huang and Oeding's \cite{HuangOeding}.
The conclusion is likewise classical as a description: \(A^2=\lambda I\)
makes \(A/\sqrt\lambda\) a hollow symmetric involution, equivalently a
constant-diagonal rank-three projection or an equal-norm tight frame of six
vectors in \(\mathbf R^3\), and the equal-modulus case is the classical
correspondence between conference matrices, regular two-graphs, and
equiangular tight frames.  What we prove is the converse direction: nonzero
proportionality of the two cubics forces the order and the scalar square,
with the order-six nondegeneracy form in
Proposition~\ref{prop:nonsingular-complementary-minors}.  We have not located
either statement in the bounded audit documented with the distributed source
archive, whose scope and access gaps are summarized in
Section~\ref{sec:verification}; to our knowledge neither appears in the
Pfaffian-commutator, compound-minor, determinantal-representation,
conference-matrix, tight-frame, or maximal-determinant design literature.
Greaves and Suda's two-valued fourth-order spectrum, used in
Section~\ref{sec:reconstruct-signing}, concerns principal \(4\times4\)
minors, not these complementary \(3\times3\) minors.

\paragraph{Rigidity of the equality.}
Let \(W\) be the fifteen-dimensional space of real symmetric order-six matrices
with zero diagonal, coordinatized by the entries \(a_{ij}\), and for a
\(3\)-subset \(S\) let \(h_S(A)\) be the coefficient of \(x_S\) in
\(\operatorname{Pf}[D_x,A]\).  Both \(h_S\) and
\(\tau_S(A)=a_{ij}a_{jk}a_{ki}\) are homogeneous cubics in the entries, so the
oriented equality locus
\[
 X_\varepsilon=\{A\in W:h_S(A)=4\varepsilon\tau_S(A)\ \text{for all twenty}\ S\},
 \qquad\varepsilon=\pm1,
\]
is a cone.  The representative \(C\) of Section~\ref{sec:orientation-source}
lies in \(X_{+1}\), and with it the whole line \(\mathbf R\,C\).  Nothing else
lies nearby.

\begin{theorem}[Rigidity at the golden representative]
\label{thm:golden-equality-rigidity}
The twenty functions \(F_S=h_S-4\tau_S\) have Jacobian of rank fourteen at
\(C\), with kernel the scaling line \(\mathbf R\,C\).  Near \(C\) the locus
\(X_{+1}\) is therefore exactly that line, and the same holds for
\(h_S+4\tau_S\) near every odd signed-permutation image of \(C\), the
representatives that reverse the orientation of the signed label space.
\end{theorem}

\begin{proof}
Euler's relation gives \(\mathrm dF_S(C)[C]=3F_S(C)=0\), so the kernel contains
\(\mathbf R\,C\) and the rank is at most fourteen.  The content is the reverse
inequality.

Each \(F_S\) is multilinear in the entries: \(\tau_S\) uses each edge inside
\(S\) once, and \(h_S\) is the complementary minor \(\det A[S^c,S]\) up to sign,
a determinant with one entry from each row and column of the cross block.  So
every partial derivative is the difference of two values of \(F_S\), at
\(a_e=1\) and at \(a_e=0\), and is an exact integer.  Only twelve of the fifteen
edges occur: for \(e=\{i,j\}\subset S=\{i,j,k\}\) one gets
\(\partial F_S/\partial a_e=-4\varepsilon a_{ik}a_{jk}\); for \(e\) joining
\(S\) to \(S^c\) the derivative is the corresponding signed \(2\times2\) cofactor
of \(A[S^c,S]\); and for \(e\) inside \(S^c\) it vanishes.

Let \(G\) be the stabilizer of \(C\) among the signed permutations
\(A\mapsto D_\epsilon P_\sigma AP_\sigma^{\mathsf T}D_\epsilon\), taken modulo
the global sign, which acts trivially on edge weights.  Triangle products are
invariant up to relabelling, since each vertex of a triangle carries two of its
edges and the two signs cancel, while \(h_S\) acquires
\(\operatorname{sgn}(\sigma)\det D_\epsilon\).  Both factors are \(+1\)
throughout \(G\), so \(G\) permutes the \(F_S\) and the Jacobian
\(J\colon W\to\mathbf R^{20}\) is \(G\)-equivariant.  Here \(G\) has order sixty,
with element orders \(1,2,3,5\) occurring \(1,15,20,24\) times, so it is the
alternating group of degree five: the orientation-preserving half of the
two-graph automorphism group recovered in the pentagon normalization above.

Hence \(\ker J\) is a \(G\)-submodule.  Every irreducible representation of that
group has a nonzero vector fixed by a subgroup of order three---the fixed-space
dimensions \((\chi(1)+2\chi(c_3))/3\) are \(1,1,1,2,1\) in dimensions
\(1,3,3,4,5\)---so if \(\ker J\) were larger than \(\mathbf R\,C\), complete
reducibility would split off a nonzero complementary submodule and give
\[
 \dim\bigl(\ker J\cap W^{\langle h\rangle}\bigr)\ge2
\]
for any \(h\in G\) of order three.  The computation therefore takes place in
\(W^{\langle h\rangle}\) rather than in \(W\).

Take \(h=(0\,2\,4)(1\,3\,5)\) with signs \((+,-,-,-,-,+)\).  Its five edge
orbits carry trivial sign holonomy, so their signed sums
\[
\begin{aligned}
 u_1&=e_{01}+e_{23}-e_{45}, &
 u_2&=e_{02}-e_{04}+e_{24}, &
 u_3&=e_{03}-e_{14}-e_{25},\\
 u_4&=e_{05}+e_{12}+e_{34}, &
 u_5&=e_{13}+e_{15}-e_{35} &&
\end{aligned}
\]
are a basis of the five-dimensional space \(W^{\langle h\rangle}\), in which
\(C=u_1+u_2+u_3-u_4-u_5\).  On this basis the twenty rows of \(J\) collapse to
eight, each divided by the content of its entries:
\begin{equation}
\label{tab:reduced-jacobian}
\begin{array}{c|rrrrr}
 \text{triples} & u_1 & u_2 & u_3 & u_4 & u_5\\\hline
 012,\,045,\,234 &  1 &  0 & -1 & -1 &  1\\
 013,\,145,\,235 &  2 & -1 &  1 &  3 & -1\\
 014,\,023,\,245 & -2 & -1 & -1 & -3 & -1\\
 015,\,123,\,345 & -1 &  1 &  1 &  1 &  0\\
 024             & -1 &  2 &  0 &  1 &  0\\
 025,\,034,\,124 &  3 & -1 & -1 &  2 & -1\\
 035,\,125,\,134 & -3 & -1 &  1 & -2 & -1\\
 135             &  1 &  0 &  0 & -1 &  2
\end{array}
\end{equation}
Every row annihilates \((1,1,1,-1,-1)\), as Euler's relation requires.  The rows
from \(012\), \(013\), \(015\) and \(024\) have determinant \(-5\) in the columns
\(u_1,\dots,u_4\), so the rank on \(W^{\langle h\rangle}\) is four and
\(\ker J\cap W^{\langle h\rangle}\) is the line through \(C\).  That contradicts
the displayed inequality unless \(\ker J=\mathbf R\,C\).

Fourteen of the twenty differentials \(\mathrm dF_S(C)\) are independent, so
by the implicit function theorem the common zero set of those fourteen
functions is a one-dimensional submanifold near \(C\).  It contains
\(X_{+1}\), which contains the scaling line; two one-dimensional
submanifolds, one inside the other through \(C\), coincide near \(C\), so all
three sets do.  An odd signed permutation negates \(h_S\) and fixes
\(\tau_S\), carrying the whole argument to \(X_{-1}\).
\end{proof}

The rigidity has a module form, which says what the equality does to the
conference directions.  Write \(T\) for the tangent space at \(C\) to the locus
\(\{A^2=\lambda I\}\).  Its dimension is forced rather than computed.  A real
symmetric hollow \(A\) satisfies \(A^2=\lambda I\) exactly when
\(A/\sqrt\lambda=2P-I\) for an orthogonal projection \(P\) of rank three with
every diagonal entry \(\tfrac12\), that is, exactly when the rows of any
factorization of \(P\) form an equal-norm tight frame of six vectors in
\(\mathbf R^3\).  Those projections are cut out of \(\mathrm{Gr}(3,\mathbf R^6)\),
of dimension nine, by the six equations \(P_{ii}=\tfrac12\), whose sum holds
identically on the Grassmannian; at most five are independent, so the projections
form a family of dimension at least four, and \(\dim T\ge5\) after the scaling
direction is added.

Equality holds, by the same mechanism that drives the recognition theorem.
Differentiating \(A^2=\lambda I\) gives \(CX+XC=\mu(X)I\) for \(X\in T\), so
\(T\) lies in the linear space
\[
 L=\{X\in W:CX+XC\in\mathbf R\,I\},
\]
on which \(\mu\) is a linear functional.  Writing \(P_\pm=(I\pm C/\sqrt5)/2\)
for the spectral projections of \(C\), so that
\(CX+XC=2\sqrt5\,(P_+XP_+-P_-XP_-)\), membership in \(L\) says
\(P_+XP_+=cP_+\) and \(P_-XP_-=-cP_-\) for a single scalar \(c\); the scalar
part is \(c(P_+-P_-)=cC/\sqrt5\), so \(L=\mathbf R\,C\oplus\ker\mu\) with
\(\ker\mu=\{X\in W:P_+XP_+=P_-XP_-=0\}\).  Forgetting hollowness, a symmetric
\(X\) with both spectral blocks zero is \(B+B^{\mathsf T}\) for an arbitrary
\(B\colon\operatorname{im}P_-\to\operatorname{im}P_+\), nine parameters, with
diagonal \(X_{ii}=2\langle p_i,Bq_i\rangle\) where \(p_i=P_+e_i\) and
\(q_i=P_-e_i\).  Hollowness pairs \(B\) against the six matrices
\(p_iq_i^{\mathsf T}\), and a relation
\(\sum_ic_ip_iq_i^{\mathsf T}=P_+D_cP_-=0\) makes the diagonal matrix \(D_c\)
commute with \(P_+\), hence with \(C\), hence forces \(c\) constant, since no
off-diagonal entry of \(C\) vanishes; as \(\sum_ip_iq_i^{\mathsf T}=P_+P_-=0\),
the relations are exactly the constants.  Five of the six conditions are
therefore independent, \(\dim\ker\mu=4\) and \(\dim L=5\), and the lower bound
gives \(T=L\).

The splitting is an eigenspace decomposition: multiplying \(CX+XC=\mu(X)I\)
by \(C\) and using \(C^2=5I\) gives
\[
 CXC=\mu(X)C-5X,
\]
so \(X\mapsto\tfrac15CXC\) is an involution of \(T\) with \(+1\)-eigenspace
\(\mathbf R\,C\) and \(-1\)-eigenspace \(\ker\mu\).  As a \(G\)-module
\(W\cong\mathbf1\oplus\mathbf4\oplus\mathbf5\oplus\mathbf5\), from the
character \((15,3,0,0,0)\) on the classes of orders \(1,2,3,5,5\); the only
four-dimensional submodule is the irreducible \(\mathbf4\), so \(T\) is
exactly \(\mathbf1\oplus\mathbf4\).  By
Theorem~\ref{thm:golden-equality-rigidity} the Jacobian \(J\) is injective on
\(\ker\mu\): the conference deformations the cubic equality kills are exactly
the scalings.  This makes the conference count and the cubic count one
statement about the four-dimensional constituent.

On the sign locus the cubic identity can be dropped altogether.

\begin{proposition}[Nonsingular complementary minors]
\label{prop:nonsingular-complementary-minors}
Let \(A\) be symmetric of order six with zero diagonal and every off-diagonal
entry \(\pm1\).  Then \(A^2=5I\) if and only if \(\det A[S^c,S]\ne0\) for every
\(3\)-subset \(S\), and in that case every one of those twenty determinants is
\(\pm4\).
\end{proposition}

\begin{proof}
A \(3\times3\) matrix with entries \(\pm1\) has determinant \(0\) or \(\pm4\):
subtracting its first row from the other two leaves entries in \(\{0,\pm2\}\)
there, so the determinant is divisible by \(4\), while Hadamard's inequality
bounds it by \(3\sqrt3<6\).  Such a matrix is singular exactly when two of its
rows are equal or opposite.  One direction is clear.  For the other, let
\(g_1,g_2,g_3\) be the pairwise inner products of the rows; if no two rows are
equal or opposite then each \(g_i=\pm1\), and the Gram determinant
\(27+2g_1g_2g_3-3(g_1^2+g_2^2+g_3^2)=18+2g_1g_2g_3\) lies in \(\{16,20\}\).  It
is the square of the determinant, hence lies in \(\{0,16\}\), so it is \(16\).

Now fix \(p\ne q\), let \(r\) be any index other than \(p\) and \(q\), and put
\(S=\{1,\dots,6\}\setminus\{p,q,r\}\), so that \(S^c=\{p,q,r\}\).  The rows of
\(A[S^c,S]\) indexed by \(p\) and \(q\) are \((a_{pm})_{m\in S}\) and
\((a_{qm})_{m\in S}\), and since the sum defining \((A^2)_{pq}\) omits
\(m=p,q\) and runs over \(S\cup\{r\}\),
\begin{equation}
\label{eq:complementary-inner-product}
 \sum_{m\in S}a_{pm}a_{qm}=(A^2)_{pq}-a_{pr}a_{qr}.
\end{equation}

Suppose every complementary minor is nonzero.  Then no two rows of
\(A[S^c,S]\) are equal or opposite, so the left side of
\eqref{eq:complementary-inner-product}, a sum of three signs, is not \(\pm3\);
hence \(\lvert(A^2)_{pq}-a_{pr}a_{qr}\rvert=1\) for each of the four admissible
\(r\).  Put \(t_r=a_{pr}a_{qr}\in\{\pm1\}\) and let \(k\) be the number of
\(r\) with \(t_r=1\), so \((A^2)_{pq}=\sum_rt_r=2k-4\).  If \(k\ge1\) some
\(t_r=1\) and \(\lvert2k-5\rvert=1\), forcing \(k\in\{2,3\}\); if \(k\le3\) some
\(t_r=-1\) and \(\lvert2k-3\rvert=1\), forcing \(k\in\{1,2\}\).  Every
\(k\in\{0,\dots,4\}\) meets at least one of the two hypotheses, and only
\(k=2\) survives both.  Hence \((A^2)_{pq}=0\), while
\((A^2)_{pp}=\sum_{m\ne p}a_{pm}^2=5\).

Conversely, if \(A^2=5I\) then \eqref{eq:complementary-inner-product} gives
\(\sum_{m\in S}a_{pm}a_{qm}=-a_{pr}a_{qr}=\pm1\) for every pair of rows of every
complementary block, so no such block has two equal or opposite rows and all
twenty determinants are \(\pm4\).
\end{proof}

This is where the magnitude of the factor \(4\) in
Theorem~\ref{thm:triangle-pfaffian-recognition} comes from: a \(3\times3\)
sign matrix has determinant \(0\) or \(\pm4\) and nothing else, so once the
complementary minors are forced away from zero the magnitude has no freedom
left.  The two orbit computations displayed earlier supply what nondegeneracy
cannot, the single sign shared by all twenty triples.

\paragraph{Why the determinant.}
The restricted determinant above is intrinsic; its scalar is fixed by the
marked determinant-line orientations.  A permanent has no analogous
basis-free meaning: changing orthonormal frames sends a matrix \(B\) to
\(R_-^{\mathsf T}BR_+\); already for \(B=I_3\), right multiplication by a
rotation through angle \(\theta\) in a coordinate plane changes its permanent
from \(1\) to \(\cos(2\theta)\).  Thus the permanent supplies no further cubic
shadow of the marked operator.

No physical or exceptional-group interpretation is used here.  The retained
chain is incidence descent, conference carrier, classical cubics, and the
harmonic return of Section~\ref{sec:harmonic}.

The next two subsections use the same conference
carrier but supply no hypothesis for Theorem~\ref{thm:harmonic-main}.  The
switching recovery inside Theorem~\ref{thm:operator-shadows} concerns the fixed
order-six conference class; the later reconstruction theorem is a global
faithfulness statement for aligned four-set data.

\subsection{Exchange spectra of symmetric conference matrices}
\label{sec:exchange-spectra}

The next theorem is a general spectral formula for every symmetric conference
order; order six is its answer, not its hypothesis.  It explains why this
six-axis carrier has no cut-independent higher-order analogue.  Let \(C=C^{\mathsf T}\) be any symmetric conference
matrix of order \(2d\), that is, a zero-diagonal matrix with off-diagonal
entries in \(\{\pm1\}\) and \(C^2=qI\), where \(q=2d-1\).  Its zero trace
gives two eigenspaces of dimension \(d\).  For a balanced half
\(Y\), order the coordinates by \((Y,Y^c)\) and write
\[
 C=\begin{pmatrix}A&R\\R^{\mathsf T}&E\end{pmatrix}.
\]
Let \(D_Y\) be the diagonal sign involution of this cut, put
\(P_\pm=(I\pm C/\sqrt q)/2\), and choose an isometry
\(Q_+:\mathbf R^d\to\operatorname{im}P_+\).  The balanced exchange operator
is
\[
 H_Y=Q_+^{\mathsf T}D_YP_-D_YQ_+.
\]

\begin{theorem}[Balanced exchange rigidity]
\label{thm:balanced-exchange-rigidity}
With the notation above,
\[
 \Spec(H_Y)=\Spec\!\left(\frac1qRR^{\mathsf T}\right)
 =\left\{1-\frac{\alpha_1^2}{q},\ldots,
          1-\frac{\alpha_d^2}{q}\right\},
\]
where \(\alpha_1,\ldots,\alpha_d\) are the eigenvalues of \(A\).  The
spectrum is independent of the balanced half \(Y\) if and only if
\(d\leq3\).  Thus order six is the unique nontrivial realized symmetric
conference order with this property, and its exchange spectrum is
\[
 \left\{\frac15,\frac45,\frac45\right\}.
\]

For a four-set \(K\), let \(w(K)\) be the sum of its three signed
Hamilton-cycle products.  Call \(K\) \emph{aligned} if \(w(K)=3\); in
standard two-graph terminology this is the union of the coherent and
incoherent four-set types.  Let \(c_Y\) count the aligned four-sets in
\(Y\).  Then
\[
 \operatorname{tr}(H_Y)=\frac{d^2}{q},\qquad
 \operatorname{tr}(H_Y^2)=\frac{F_d+32c_Y}{q^2},
\]
where
\[
 F_d=dq^2-2qd(d-1)+d(d-1)+12\binom d3-8\binom d4.
\]
Hence the first exchange moment is universal, whereas the second is
cut-independent exactly in the stated small orders.
\end{theorem}

\begin{proof}
Put \(Q=C/\sqrt q\) and \(L=[D_Y,Q]\).  The commutator anticommutes with
\(Q\), and on its positive eigenspace one has \(Lv=2P_-D_Yv\).  In cut
coordinates,
\[
 L=\frac2{\sqrt q}
 \begin{pmatrix}0&R\\-R^{\mathsf T}&0\end{pmatrix}.
\]
Consequently \(4H_Y\) is \(L^*L\) restricted to the positive eigenspace.
Thus \(H_Y\) has the squared singular values of \(R/\sqrt q\).  The
\(Y\)-principal block of \(C^2=qI\) is
\[
 RR^{\mathsf T}=qI-A^2,
\]
which proves the spectral formula.  Complementary principal-block spectral
relations of this kind are classical \cite{HaemersParsaeiMajd}.

Since \(A\) is a zero-diagonal sign matrix,
\(\operatorname{tr}(A^2)=d(d-1)\), giving the first trace formula.  Sorting
the closed four-walks by support gives
\[
 \operatorname{tr}(A^4)=d(d-1)+12\binom d3
   +8\sum_{\substack{K\subset Y\\|K|=4}}w(K).
\]
Every edge occurs twice in the product of the three cycle signs, so
\(w(K)\in\{3,-1\}\).  Substitution yields the displayed formula for
\(\operatorname{tr}(H_Y^2)\).  This is the exact finite version of the
closed-walk organization used for conference-frame subensembles in
\cite{MagsinoMixonParshall}.

Suppose now that \(d\geq4\) and even the second exchange moment is independent
of \(Y\).  Then the sums of the aligned-four-set indicator over all balanced
halves are equal, and the indicator is therefore constant on four-sets.  Only
this consequence of the rank formula for inclusion matrices
\cite[Theorem~1]{JolliffeInclusion} is needed, and it follows from a
one-element exchange: two balanced halves differing in \(a\) versus \(b\) have
equal sums, so
\[
 \sum_{\substack{S\subset Y\cap Y'\\ |S|=3}}
   \bigl(f(S\cup\{a\})-f(S\cup\{b\})\bigr)=0
\]
for the aligned indicator \(f\).  That is the same hypothesis one subset size
lower, for the difference function on the remaining \(2d-2\) points.  Two
more exchanges of the same kind leave a function of single points whose sum
over every \((d-3)\)-subset vanishes; comparing two such subsets differing in
one element makes it constant, hence zero.  Now climb back up: a function
whose one-element differences all vanish is constant, since one-element
exchanges connect the subsets of a fixed size, and a constant with vanishing
sum over a nonempty family is zero.  Applied to the second and then the first
difference this makes both vanish identically, and vanishing first
differences make \(f\) constant on four-sets.  Every subset and exchange used
exists because \(d\geq4\).

Let \(w\) be the common value of \(w(K)\) and apply the closed-walk count to
the whole matrix rather than to a half.  Reading the diagonal of \(C^2=qI\)
gives \(q=2d-1\), so \(\operatorname{tr}(C^4)=2d(2d-1)^2\), while the count
gives
\[
 \operatorname{tr}(C^4)=2d(2d-1)+12\binom{2d}3+8w\binom{2d}4 .
\]
Since \(2d(2d-1)^2-2d(2d-1)-12\binom{2d}3=-2d(2d-1)(2d-2)\) and
\(8\binom{2d}4=2d(2d-1)(2d-2)(2d-3)/3\), dividing by
\(2d(2d-1)(2d-2)\neq0\) leaves
\[
 (2d-3)\,w=-3 .
\]
Hence \(w=3\) forces \(2d=2\) and \(w=-1\) forces \(2d=6\), and neither is
compatible with \(d\geq4\).  At the surviving order the same equation says
that every four-set of the order-six conference matrix carries weight \(-1\),
so that matrix has no aligned four-set.

For \(d=1\) the conclusion is immediate, and for \(d=2\) every possible
two-by-two principal sign block squares to \(I\).  For \(d=3\), if \(\tau\)
is the product of the three edge signs, then
\[
 A^2=2I+\tau A,
\]
and \(\det(tI-A)=t^3-3t-2\tau\).  This polynomial factors as
\((t-2)(t+1)^2\) or \((t+2)(t-1)^2\), so \(A^2\) always has spectrum
\(\{4,1,1\}\).  The spectral formula then gives
\(\{1/5,4/5,4/5\}\).
\end{proof}

The obstruction is local, but the aligned sets retain uniform aggregate
statistics.  Greaves and Suda proved that the determinant-\((-3)\)
four-subsets of a symmetric conference matrix form a
\(3\)-\((2d,4,(d-3)/2)\) design \cite[Example~2.3]{GreavesSuda}; the identity
\(\det C[K]=3-2w(K)\) identifies their blocks with the aligned sets above.
The regular-two-graph parameters give the same identification
\cite[Section~2.2 and Proposition~4.1]{GillespieTwoGraph}.  Classical
Johnson-scheme inclusion calculus \cite[Section~8]{IntersectionMatrices}
therefore gives, for a uniformly chosen balanced half and \(d\geq4\),
\[
 \mathbf E c_Y=\rho\binom d4,\qquad
 \operatorname{Var}(c_Y)=
 \frac{\binom{2d}{4}\binom{2d-8}{d-4}}{\binom{2d}{d}}
 \rho(1-\rho),
 \qquad \rho=\frac{d-3}{2(2d-3)}.
\]
These formulas also hold modulo complementation because \(c_Y=c_{Y^c}\).
At order ten they give \(\mathbf E c_Y=5/7\) and
\(\operatorname{Var}(c_Y)=10/49\), hence
\(\mathbf E[c_Y(c_Y-1)]=0\).  Since \(c_Y\) is a nonnegative integer, it is
therefore \(0\) or \(1\) on every cut.  The 126 projective cuts consequently split
into 36 with \(c_Y=0\) and 90 with \(c_Y=1\).  The known design fixes this
second-moment split without making the complete higher-order spectrum uniform.

\subsection{Reconstructing the signing}
\label{sec:reconstruct-signing}
The design still retains the full discrete signing, despite fixing only
aggregate low moments.  Let
\(\tau:\binom V3\to\mathbf F_2\) be a two-graph on a finite set \(V\), meaning
that it satisfies \eqref{eq:two-graph-parity} with \(X\) replaced by \(V\).
Let \(\mathcal A(\tau)\) be the four-sets on which the four values of
\(\tau\) are all equal.  Complementing \(\tau\) preserves this family.

Equation~\eqref{eq:two-graph-parity} is the definition of a two-graph, and the
correspondence between two-graphs and switching classes of graphs, used below
as the two-graph equation, is likewise standard: two-graphs are due to Higman,
with Taylor the standard reference and Seidel's switching the graph-side
description; see Brouwer and Van Maldeghem \cite[\S1.1.12]{BrouwerVanMaldeghem}
for both, with attributions.

\begin{theorem}[Aligned-design faithfulness]
\label{thm:aligned-faithfulness}
If \(|V|\geq7\), then \(\mathcal A(\tau)\) determines \(\tau\) up to
complement, and seven is the least bound with this property.

For a symmetric conference matrix of order \(n\geq10\), its marked family of
principal four-subsets with determinant \(-3\), which is the design of
Greaves and Suda \cite[Table~1 and Example~2.3]{GreavesSuda}, therefore
determines its signing up to diagonal switching and global negation.  After
one aligned four-set \(Q\) is known, the
\[
 1+4(n-4)+6\binom{n-4}{2}=3n^2-23n+45
\]
tests on four-sets meeting \(Q\) in at least two points suffice, and the
decoder in the proof uses exactly this many.  An aligned anchor is found
deterministically with at most twenty tests.  Each of those queries reads one
bit, whether a single fourth-order principal minor equals \(-3\); the
reconstruction consumes no minor value and no minor of any other order.
One calibrated triangle product selects between the two global orientations.
\end{theorem}

\begin{proof}
Fix \(r\in V\) and form a graph \(G_r\) on \(V\setminus\{r\}\), with \(ij\)
an edge when \(\tau(rij)=1\).  The two-graph equation gives
\[
 \tau(ijk)=\tau(rij)+\tau(rik)+\tau(rjk).
\]
Thus \(\{r,i,j,k\}\) is aligned exactly when \(ijk\) is a clique or an
independent triple of \(G_r\).  The equality \(R(3,3)=6\) supplies an aligned
anchor on every seven-set.

It remains to prove faithfulness on seven points.  Fix an aligned
\(Q=\{1,2,3,4\}\), complement one candidate if necessary, and switch graph
representatives so that both vanish on \(Q\).  Switch each outside vertex
once more so that its edge to vertex \(4\) is zero.  An outside point \(x\)
then determines a uniquely normalized cut
\[
 p_x=(e_{1x},e_{2x},e_{3x},e_{4x})
 \in\mathbf F_2^4,\qquad (p_x)_4=0.
\]
The four tests meeting \(Q\) in three points have the following complete
signature table; an entry \(ijk\) means that \(\{i,j,k,x\}\) is aligned.
\[
\begin{array}{c|c@{\qquad}c|c}
p_x&\text{aligned triples}&p_x&\text{aligned triples}\\ \hline
0000&123,124,134,234&1000&234\\
0100&134&0010&124\\
1110&123&1100,1010,0110&\varnothing
\end{array}
\]
Thus they recover \(p_x\), except that the three balanced cuts
\(B_{12}=1100\), \(B_{13}=1010\), and \(B_{14}=0110\) have the same empty
signature.

We record the remaining finite step explicitly.  For outside points \(x,y\),
put \(p=p_x\), \(s=p_y\), \(e=e_{xy}\), and let
\(\mathcal B(p,s,e)\) be the set of pairs \(ij\subset Q\) for which
\(\{i,j,x,y\}\) is aligned.  Its four triangle values give
\begin{equation}
 \label{eq:aligned-pair-criterion}
 ij\in\mathcal B(p,s,e)
 \quad\Longleftrightarrow\quad
 p_i+s_i=p_j+s_j\ \text{ and }\ e=p_j+s_i.
\end{equation}
This criterion gives a complete classification.  If \(p\) and \(s\) are
already known, some pair has \(p_i+s_i=p_j+s_j\), and exactly one value of
\(e\) makes that pair aligned.  If exactly one cut is balanced, relabelling
\(Q\) reduces the known cut to \(0000\) or \(1000\); substituting the three
balanced cuts in \eqref{eq:aligned-pair-criterion} gives
\[
\begin{array}{c|c|c|c}
p&s&\mathcal B(p,s,0)&\mathcal B(p,s,1)\\ \hline
0000&B_{12}&34&12\\
0000&B_{13}&24&13\\
0000&B_{14}&14&23\\
1000&B_{12}&34&13,14\\
1000&B_{13}&24&12,14\\
1000&B_{14}&12,13&23
\end{array}
\]
so both the balanced cut and \(e\) are recovered.  When both cuts are
balanced, the remaining table is
\[
\begin{array}{c|c|c}
\{p,s\}&\mathcal B(p,s,0)&\mathcal B(p,s,1)\\ \hline
\{B_{12},B_{12}\}&12,34&13,14,23,24\\
\{B_{13},B_{13}\}&13,24&12,14,23,34\\
\{B_{14},B_{14}\}&14,23&12,13,24,34\\
\{B_{12},B_{13}\}&23&14\\
\{B_{12},B_{14}\}&13&24\\
\{B_{13},B_{14}\}&12&34
\end{array}
\]
The rows are distinct.  The sole loss of information is the order of two
distinct balanced cuts: interchanging \(p\) and \(s\) preserves the six pair
tests and leaves \(e\) fixed.

There are three outside points.  Suppose one cut changes between the two
candidates.  For each of the other two points the classification forces the
pair to consist of distinct balanced cuts and swaps them.  Those two other
cuts are therefore equal to the same changed cut.  Their mutual pair cannot
be a distinct-cut swap, so it must agree pointwise, contradicting either of
the first two swaps.  No cut changes, and then
\eqref{eq:aligned-pair-criterion} recovers every outside
edge.  The candidates agree after the initial possible complement.

For larger \(V\), apply the seven-point result to every seven-set.  Adjacent
seven-sets in the connected Johnson graph share six vertices.  On any triple
in that intersection the two restrictions cannot choose opposite complement
conventions, so their complement choices agree.  This proves faithfulness.
This proof is also the promised decoder.  Fix a root and six other vertices
and test the twenty four-sets formed by the root and a triple of the six;
\(R(3,3)=6\) supplies an aligned anchor \(Q\).  Query all four-sets meeting
\(Q\) in at least two points.  First reconstruct one seven-set
\(Q\cup\{x,y,z\}\).  For each remaining vertex \(u\), apply the same
seven-point tables to \(Q\cup\{x,y,u\}\); agreement on the fixed six-set
\(Q\cup\{x,y\}\) selects the same complement convention and recovers the
cut of \(u\).  Once all cuts are known,
\eqref{eq:aligned-pair-criterion} recovers every outside edge.
No further query is used, so the number of selected tests is exactly the
quadratic count in the theorem.

For a conference signing, the identity
\(\det C[Q]=3-2w(Q)\) identifies the determinant-\((-3)\) family with
\(\mathcal A(\tau)\).  Choosing a root and switching its incident signs to
one recovers every matrix entry from the rooted triangle values.  The sole
remaining complement sends \(C\) to \(-C\), and one triangle product fixes
that orientation.  Minimality of the bound seven is the six-point pair of
Remark~\ref{rem:six-point-witness}.
\end{proof}

\begin{remark}[Seven points is sharp]
\label{rem:six-point-witness}
Six points do not suffice, and not only because a two-graph may have no
aligned four-set at all.  Represent a two-graph on \(\{0,\dots,5\}\) by the
unique graph in its switching class in which \(0\) is isolated, and take
\[
 G=\{12,\,15,\,24,\,25,\,35\},
 \qquad
 H=\{13,\,14,\,24,\,34,\,35\}.
\]
Both two-graphs have aligned family exactly
\(\bigl\{\{0,1,2,5\},\{0,1,3,4\}\bigr\}\), and \(H\) is neither \(G\) nor its
complement.  So \(\mathcal A\) does not determine a two-graph on six points up
to complement, and the hypothesis \(|V|\ge7\) of
Theorem~\ref{thm:aligned-faithfulness} cannot be weakened.
\end{remark}

\begin{remark}[What the query count is measured against]
\label{rem:alignment-query-count}
For a general symmetric matrix, Holtz and Sturmfels identify the
principal-minor ambiguity as diagonal sign similarity under their strict
nondegeneracy hypothesis \cite[Thm.~6]{HoltzSturmfels}.  No such hypothesis
is needed here: for a Seidel matrix the order-three minors are twice its
triangle signs, so they recover the switching class directly.  Against an
oracle
returning minor \emph{values} the problem is solved in \(O(n^2)\) queries and
that count is asymptotically optimal
\cite{RisingKuleszaTaskar,BrunelUrschel}; for a Seidel matrix the cycle-basis
algorithm needs only the \(\binom n2-n+1\) principal minors of order three
\cite{BrunelSignedDPP}, each of which is twice a triangle sign and so returns
the two-graph outright.  The count above improves on none of that and is not
offered as an algorithm.  It is what a decoder pays when it is given only the
fourth-order indicator, which is invariant under global negation as the
order-three data is not.

Two lower bounds hold for that data.  Two-graphs on \(n\) points form an
\(\mathbf F_2\)-space of dimension \(\binom n2-n+1\); each test returns one
bit and determination up to complement allows fibres of size two, so any
family of tests --- and any decision tree, by its leaf count --- has at least
\(\binom n2-n\) members.  Families fixed in advance obey a stronger bound,
because the answers are biased.  A uniformly random two-graph restricts to a
uniformly random two-graph on any four points, and two of those eight are
aligned, so every alignment test answers yes with probability exactly
\(1/4\).  The information-theoretic bound of combinatorial search theory,
subadditivity of entropy applied to the entropy \(\binom{n-1}{2}-1\) of the
complement pair, which is what the answers determine, gives
\[
 k\;\ge\;\frac{\binom{n-1}{2}-1}{H(1/4)}
 \;>\;1.2326\Bigl(\tbinom{n-1}{2}-1\Bigr)\;\approx\;0.616\,n^2,
\]
where \(H(p)=-p\log_2p-(1-p)\log_2(1-p)\), for every family of \(k\) tests
that determines the two-graph up to complement.  The ratio of the exhibited
count \(3n^2-23n+45\) to this floor tends to \(4.87\), and the constant
\(3\) is not known to be optimal.

The floor does not apply to a decoder that chooses each test in the light of
earlier answers, and for large \(n\) such a decoder does better than any
fixed family can, using \(\binom n2+O(n)\) tests --- the counting bound to
leading order.  Read the two-graph on seven points, which costs at most the
\(\binom74\) tests available there, and then add one point at a time.
Attaching a point \(v\) to six points carrying a known restriction costs a
bounded number of tests, because the twenty tests meeting \(v\) inside those
seven points determine its five rooted values: the fifteen tests avoiding
\(v\) are already known, and the theorem separates on seven points.  Every
further rooted value at \(v\) then costs a single test.  Indeed, for a new
old vertex \(u\), the two hypothetical values of \(\tau(ruv)\), together
with the already recovered data, define two candidate restrictions on any
seven-set containing \(u,v\) and five of the six anchor vertices.  They agree
on the anchor five-set and hence are not complements.  By the seven-point
theorem some alignment test distinguishes them.  Every test not containing
both \(u\) and \(v\) is already decided, so one of the tests through
\(u,v\) and two anchors distinguishes the two values.  The decoder evaluates
both hypothetical outcomes, selects such a test, and its one-bit answer fixes
\(\tau(ruv)\).  So the price of the coherence restriction is a
price of fixing the tests in advance: adaptively the fourth-order indicator
matches the order-three minor values to leading order, both at \(n^2/2\).

Tests indexed by four-sets occur elsewhere as a primitive observation.  The
causal-clustering algorithm of Kummerfeld and Ramsey \cite{KummerfeldRamsey}
decides one coherence bit per quartet of observed variables, up to
\(\binom n4\) of them, each by a pair of statistical tests on sampled data.
That bit reports whether covariance minors vanish under a one-factor
measurement model, not whether four triple signs cohere, so the two counts
measure different work and are not compared here.
\end{remark}

\smallskip
\noindent\emph{The same statement in principal minors.}\par
For a Seidel matrix the theorem is a statement about principal minors, and it
is worth recording which ones it uses.  The principal minors of order at most three trivialize the
problem: the \(1\times1\) minors vanish, the \(2\times2\) minors are all
\(-1\), and the third-order minors return the two-graph outright, as
Remark~\ref{rem:alignment-query-count} records.  The theorem consults none of
them.  Its data are the fourth-order principal minors, through
\(\det C[Q]=3-2w(Q)\) below; that these take only the two values \(-3\) and
\(5\) is Greaves and Suda's Table~1, and the \(-3\) fibre is their design.
They are invariant under global negation, so they cannot distinguish \(C\)
from \(-C\) and see each four-set only through the single bit
\(\det C[Q]\in\{-3,5\}\).  The content is that this negation-invariant
fourth-order data still determines the matrix up to exactly that ambiguity,
that seven vertices are enough and six are not by
Remark~\ref{rem:six-point-witness}, and that a quadratic family of these
determinants suffices.

This theorem belongs to the reconstruction-up-to-complementation literature.
On four vertices a two-graph has zero, two, or four triples, so equality of
aligned families is precisely four-hypomorphy up to complementation within the
class of two-graphs.  Dammak, Lopez, Pouzet, and Si Kaddour prove the analogous
four-local reconstruction theorem for ordinary graphs
\cite{DammakGraphHypomorphy}.  For arbitrary \(3\)-uniform hypergraphs,
Pouzet and Si Kaddour prove that the least local size yielding eventual
reconstruction is five \cite{PouzetHypergraphIsomorphy}.  The two-graph parity
axiom therefore lowers that general threshold to four, with ambient order
seven sharp.  That sharpness is a separate statement from Dammak, Lopez,
Pouzet, and Si Kaddour's: their \(v\ge7\) is the endpoint of the admissible
range \(4\le k\le v-3\), and the sharpness they remark on is in \(k\), so no
six-point failure follows from their work for either observable.
The homogeneous-triple analysis of ordinary graphs
\cite{PouzetHomogeneousTriples} supplies the one-anchor viewpoint; the
away-from-anchor tests above remove its balanced-cut ambiguity.

We have not located this sharp two-graph specialization in the bounded audit
documented with the distributed source archive; the paper claims no priority
for the general hypomorphy framework or the homogeneous-triple method.  The
scope and access gaps of that audit are summarized in
Section~\ref{sec:verification}.

The neighboring framework is spectral monomorphy of principal submatrices
\cite{AttasSpectralMonomorphy,BoussairiTwoGraphs}.  Here only the squared
spectrum is assumed constant.  We have not located the balanced
singular-spectral classification above in that literature; the block identity,
closed-walk expansion, inclusion rank, and aligned design are credited as
indicated.  The bounded search scope is recorded with the distributed source
archive and summarized in Section~\ref{sec:verification}.

This exchange invariant deliberately forgets signs: it sees \(A^2\), or
equivalently \(RR^{\mathsf T}\).  The cubic stage above retains the
determinant-line orientation.  Cut-independent squared singular spectrum,
determinant orientation, and the later cubic shadows are therefore distinct
claims.
